``Cái hộp ấy à?'' Tôi lên tiếng hỏi cô bạn đang hớn hở trước mặt trong
kh dọn lại tập sách bỏ vào cặp. ``Cậu vẫn còn nhớ sao Kuro?''

``Nhớ chứ! Đã một khoảng thời gian khá lâu kể từ lúc hai đứa mình đem
chôn nó rồi mà!'' Cậu ấy hí hửng trả lời tôi. ``Tớ thật sự rất tò mò
rằng hồi đó chúng ta bỏ gì vào trong đó!''

Trông Kuro có vẻ hào hứng dữ nhỉ? Mà, chính tôi cũng đã quên mất sự tồn
tại của cái hộp đó rồi, mãi đến khi cậu ấy nhắc lại thì tôi mới sực nhớ
ra rằng đúng là hai đứa bọn tôi hồi nhỏ có đem mấy cái đồ vật linh tinh
bỏ vào một cái hộp, sau đó đem chôn dưới một gốc cây, còn gốc cây đó ở
đâu thì tôi chịu.

``Thôi được rồi.'' Tuy vậy nhưng tôi vẫn không thể chiến thắng trước vẻ
mặt đang hớn ha hớn hở của Kuro nên đành thở dài đồng ý. ``Nhưng nói
trước là tớ không nhớ đã chôn cái hộp ở nơi nào đâu đấy nhé.''

``Không phải lo! Tớ biết rõ cậu não cá vàng mà Koko!'' Kuro vỗ ngực đầy
tự hào, như muốn nói rằng cứ để cậu ấy lo. ``Tớ nhớ rất rõ địa điểm hai
ta đã cùng chôn nó đấy!''

``Ai là não cá vàng cơ...'' Tôi thở dài và xách cặp đứng lên, đi theo
Kuro đang hào hứng không khác gì một đứa trẻ trước khi đi chơi.

Tôi là Kokoro, và cô bạn đang chạy tung tăng trước mắt tôi là Kurosaki,
hai đứa bọn tôi đã chơi cùng nhau kể từ thời còn bé tí, thậm chí là
trước cả khi vào mẫu giáo thì cả hai đã chơi cùng nhau rồi.

Lý do đưa cả hai bọn tôi ở bên nhau lâu đến như vậy khá là đặc biệt. Mẹ
của tôi đã ly hôn với bố khi tôi vẫn còn rất nhỏ và chuyển đến sống với
mẹ của Kuro, người cũng đã ly hôn với chồng của mình, vì tuổi tác và
tính cách khá hợp nhau nên tôi với Kuro nhanh chóng trở nên thân thiết.
Và ờm, mẹ tôi với mẹ cậu ấy sau đó công khai kết hôn với nhau nên là bên
cạnh bạn thuở nhỏ thì tôi với Kuro còn là chị em kế nữa.

Và thế là hai đứa tôi cứ thế ở bên nhau cho đến tận bây giờ, tức là cấp
ba, hai người mẹ của chúng tôi thì cứ gọi đùa bọn tôi là định mệnh của
nhau gì gì đấy, thậm chí cả hiện tại thì cũng có một vài tin đồn ở
trường cho rằng hai đứa đang hẹn hò với nhau.

Tính cách của Kuro thì vẫn vậy, vẫn cứ nhoi nhoi và thừa năng lượng đến
mức khó tin. Nhưng bằng một cách nào đó thì điểm số của cậu ta vẫn cứ
giữ được ở mức ổn định dù chả bao giờ chịu ngồi vào bàn học và luôn làm
phiền tôi khi đang ôn bài.

Quay lại vấn đề chính, khi hai đứa chuẩn bị vào lớp 1 thì Kuro đã đề
nghị là lấy những món đồ vật chứa nhiều kỷ niệm nhất của cả hai bỏ vào
một cái hộp và đem chôn và mười năm sau sẽ quay lại để đào nó lên, tôi
vẫn nhớ rõ lúc đó bản thân đã hí hửng đồng ý với chuyện này. Và kết quả
là quên khuấy mất tận mười năm.

Tự nhiên cảm thấy có lỗi ghê, chắc lát phải đãi cậu ấy một bữa vậy.

Sau một hồi đi bộ thì bọn tôi đã leo lên ngọn núi nhỏ sau trường, chỗ
này khá phổ biến cho mấy đứa con nít vào chơi bời vì ở đây có khá nhiều
cây xanh, gió mát và đường đi cũng rất an toàn, hồi nhỏ tôi với Kuro
cũng thường ra đây chơi nữa.

Kuro dắt tôi qua những nơi rất quen thuộc, những kỷ niệm mà tôi từng
quên mất bây giờ lại hiện rõ mồn một trong đầu. Và sau cùng thì đã đến
một nơi khá là kín đáo, chỗ này có một cái cây bị khoét một cái lỗ to
tướng mà thậm chí tôi bây giờ cũng có thể chui vừa vào trong đó.

``Chỗ này là...'' Tôi ngó nghiêng xung quanh. ``Căn cứ bí mật của tụi
mình nhỉ?''

``Đúng!.'' Kuro hớn hở chạy đến cái hốc mà hai đứa bọn tôi ngày xưa hay
ngồi vào trong để chơi đùa. ``Và đây cũng là nơi chôn cái hộp!''

``Ừm, tớ cũng vừa nhớ ra.''

Sau đó Kuro và tôi cùng đào lên, nơi chúng tôi dùng để chôn cái hộp là ở
trong hốc cây, nơi mà chính tôi đã gợi ý với lý do là cho dễ tìm. Tôi
cũng biết ơn bản thân khi đó đã đưa ra lựa chọn này, chứ mà để tìm được
cái hộp bé xíu đấy ở một nơi rộng mênh mông thế này thì cũng mệt mỏi ra
trò đấy.

``Thấy rồi!'' Kuro giơ cao cái hộp và hét lên như thể bản thân vừa đạt
được một thành tích cực kì đáng tự hào.

Đó là một cái hộp thiếc nhỏ không có họa tiết gì quá đặc biệt. Theo trí
nhớ của tôi thì đó là cái hộp đựng kẹo mà Kuro đã từng chia sẻ với tôi
hồi mới chuyển đến, tuy cũng chỉ là kẹo rẻ tiền thôi nhưng đó cũng là
một dấu mốc quan trọng đánh dấu cho ngày đầu tiên cả hai chúng tôi gặp
nhau.

Sau đó thì cả hai mở hộp ra, bên trong đó thì có những món đồ chơi và
trang sức cũ kĩ, tuy nhiên thì thứ tôi để ý hơn cả là cuốn nhật ký ghi
tên ``Kokoro'' trên đó, đây là nhật ký của tôi hồi đó thì phải, tôi có
bỏ vào đây à?

``A! Cái găng tay mà hồi trước cậu may tặng tớ lúc sinh nhật bốn tuổi
này!'' Kuro lấy một cái găng tay bé xíu bên trong rồi kêu lên đầy mừng
rỡ. ``Lúc cậu chìa hộp quà với cái tay đây vết băng cá nhân, tớ thật sự
đã rất lo lắng đó!''

``Ờ, tớ biết.'' Tôi nhận lấy cái găng tay và mỉm cười. ``Lúc đó cậu nắm
tay tớ khóc nức nở luôn mà.''

Sau đó cả hai đứa chúng tôi cùng nhau lấy mấy thứ bên trong ra. Những
thứ còn lại gồm có một bức tranh do Kuro vẽ bằng bút màu, bên trong là
hình ảnh tôi và cậu ấy đang nắm tay nhau. Một con búp bê mà cả hai cùng
tự dùng tiền của bản thân mua về, và một vài những thứ linh tinh khác.

``Cuối cùng là cái này nhỉ?'' Tôi cầm quyển nhật ký bên trong lên, xét
về kích thước thì nó là thứ to nhất trong hộp. Tôi định mở ra đọc nhưng
nghĩ đi nghĩ lại một hồi thì quyết định đưa sang Kuro. ``Cậu đọc đi, tớ
thấy ngại ngại sao ấy.''

``Ngại sao? Đáng yêu ghê đó.'' Kuro nở một nụ cười trông khá là đáng lo
ngại và xoa đầu tôi, đồng thời cũng nhận lấy cuốn nhật ký. ``Để xem bé
Koko lúc đó ghi gì trong nhật ký nào.''

\emph{Gửi bản thân tớ mười năm sau!}

\emph{Nếu đọc được dòng này, hẳn là đã mười năm trôi qua, và cậu với
Kuro đã đào cái hộp này lên rồi nhỉ? Không biết tớ với Kuro mười năm sau
trông như nào ha? Tớ tò mò lắm đó!}

\emph{Không biết cậu với Kuro còn chơi chung không, chứ hiện tại thì bọn
tớ vẫn đang rất thân thiết với nhau. Sẽ rất buồn nếu đọc những dòng này
một mình nên tớ mong là cậu với cậu ấy sẽ vẫn là hai người bạn thân
thiết nhất với nhau!}

\emph{À khoan, nếu được thì tớ không muốn cậu với Kuro vẫn là bạn đâu.}

\emph{Vì tớ yêu cậu ấy mà!}

\emph{Tớ hi vọng là cậu vẫn còn giữ thứ cảm xúc ấy bên trong, và tớ cũng
hi vọng là hai người sẽ thành đôi, nếu thế thì tớ sẽ có thể an tâm về
tương lai của mình rồi!}

\emph{Điều cuối cùng, chúc cậu và Kuro hạnh phúc bên nhau nhé!}

\emph{Ký tên.}

\emph{Kokoro lúc năm tuổi.}

Tôi vội vã giật lấy cuốn nhật ký từ tay của Kuro đang đỏ bừng mặt vì
những lời nói bên trong quyển nhật ký.

Thôi toang! Bây giờ tôi mới nhớ ra rằng khi đó bản thân đã quyết định
viết một bức thư gửi cho Kokoro mười năm sau, tức là chính tôi hiện tại
đây! Và với những cảm xúc bé thơ khi ấy thì tôi đã không kiềm được mà
hỏi luôn rằng tôi và Kuro đã thành đôi hay chưa.

Mà hiện tại đúng là tôi đang thích cậu ấy thật mới chết! Thế khác nào
một cách tỏ tình gián tiếp đâu! Trời ơi!

``Ừm thì...'' Kuro nghịch tóc của mìn như một cách để che đi sự xấu hổ,
đồng thời quay sang tôi, ngập ngừng cất tiếng. ``Những thứ trong nhật ký
đều là thật nhỉ?''

``Ừm.''

``Vậy hiện tại... cậu vẫn còn thích tớ chứ?''

``...Vẫn còn.''

AAAAAA! Tôi muốn tìm một cái hố để chui đầu vào quá đi! Cả hai đứa đều
là con gái, mà Kuro chắc cũng chỉ xem tôi là một người bạn thuở nhỏ siêu
thân thiết mà thôi! Xui rủi thì hai đứa chúng tôi sẽ không thể làm bạn
nữa mất!

``...Nếu tớ nói rằng, hai ta có chung một cảm xúc...'' Kuro bất chợt
tiến đến nắm lấy tay tôi. ``Liệu cậu sẽ tin chứ?''

``Ế''

Tiếng kêu ngạc nhiên chỉ vừa thoát ra thì ngay sau đó, trên môi tôi cảm
nhận được một sự mềm mại và ẩm ướt đến từ một bờ môi khác, khuôn mặt của
Kuro lúc này đang ở rất gần tôi, hai tay của cậu ấy ôm lấy eo tôi và kéo
sát người tôi vào.

``Đó là câu trả lời của tớ.'' Kuro bỏ môi tôi ra và ngại ngùng nói. ``Từ
bây giờ, mong cậu giúp đỡ nhé.''

``Ừm, tớ cũng vậy...'' Tôi gục mặt vào ngực của cậu ấy, lí nhí lên
tiếng.

Gửi bản thân tớ mười năm trước.

Hiện giờ, tớ và Kuro vẫn đang ở bên nhau. Và nhờ có bức thư mà cậu để
lại, hai đứa bọn tớ cuối cùng đã có thể tiến xa hơn trong mối quan hệ
này. Thế nên tớ xin gửi lời cảm ơn chân thành nhất đến cậu.

Điều cuối cùng, mong cậu và Kuro sẽ tiếp tục sống hạnh phúc bên nhau
trong những tháng ngày trẻ thơ và yên bình ấy.

Ký tên.

Kokoro lúc mười lăm tuổi.
